\documentclass[9pt,a4paper]{extarticle}
\usepackage[margin=12mm]{geometry}
\usepackage[T1]{fontenc}
\usepackage[scaled=0.95]{helvet}
\renewcommand{\familydefault}{\sfdefault}
\usepackage{enumitem}
\usepackage{tabularx}
\usepackage{array}
\usepackage[dvipsnames]{xcolor}
\usepackage{titlesec}
\usepackage[hidelinks]{hyperref}
\setlist{nosep,leftmargin=12pt,itemsep=1pt}
\setlength{\parindent}{0pt}
\setlength{\parskip}{2pt}
\definecolor{hd}{HTML}{1F4E79}
\titleformat{\section}{\bfseries\color{hd}\normalsize}{\thesection.}{4pt}{}[\vspace{-3pt}\color{hd!30}\hrule]
\titlespacing*{\section}{0pt}{5pt}{2pt}
\renewcommand{\arraystretch}{1.12}
\newcolumntype{Y}{>{\raggedright\arraybackslash}X}
\pagestyle{empty}

\begin{document}

{\Large\bfseries AI Lawyer India: Solution Doc (Problem 2)}\\[1pt]
{\footnotesize\color{gray}[Name] $\cdot$ [Roll no.] $\cdot$ [Branch] \quad|\quad Prototype: \textbf{\url{https://ai-lawyer-eta.vercel.app}} \quad|\quad Raw outputs: \texttt{outputs/raw/} (17 cases) \quad|\quad 8 Nov 2026}

\section{Problem choice and scope}
I picked Problem 2 because everyday legal problems in India (money not returned, a scam, a notice nobody explains) hurt many people and need no show-specific context. I went deep on \textbf{money disputes} and said what is not covered. \textbf{Covered, with verified law:} online refunds and non-delivery, e-commerce platforms, payment scams and unauthorised transactions, rental deposits, unpaid salary, private loans where documents are held, and demand notices and cheque bounce. \textbf{Not covered:} family, property title, tax, criminal defence. For those the tool still gives practical steps but labels them ``(General)'' and does not cite law. It never stops at ``see a lawyer''.

\section{Assumptions (the brief was silent)}
\begin{itemize}
\item Users are ordinary people on a phone, often not writing English; text or voice, replies in their language (script detected in code).
\item ``Today'' is the server date; criminal law means the \textbf{BNS 2023}, not the IPC. Consumer forum limits use the amount \emph{paid} (District up to Rs 50 lakh).
\item indiacode.nic.in blocked my fetches, so \textbf{33 of 53 corpus entries come from secondary sources} (bare-act sites, law-firm notes); 20 are official text. Each entry carries \texttt{source\_type} and \texttt{verified\_on}.
\item Free tiers only: Gemini free API, Vercel Hobby. No user data is stored; the case lives in the browser and is sent with each request.
\end{itemize}

\section{How it works: what AI does, what code does}
\textbf{Flow:} (1) tell your story, (2) it restates a structured case file for you to confirm, (3) it explains where you stand, cites laws and asks what proof you have, (4) it gives a step-by-step plan, the right forum, a draft notice, complaint or reply, and a verdict on whether this tool can handle the case, (5) follow-up questions, plus \textbf{Practice mode}: the AI plays the other side's lawyer, briefs you on likely tricks, asks hard questions and scores each answer.

\smallskip
\begin{tabularx}{\linewidth}{|>{\bfseries\raggedright\arraybackslash}p{26mm}|Y|}
\hline
Gemini (3.5 Flash-Lite) & Turns the messy story into facts and a timeline; explains in plain language; picks relevant provisions \emph{from the file}; drafts documents; plays the opposing lawyer. \\ \hline
Plain code & Stage order; the law file; \textbf{citation validation}; flagging of unknown section numbers, ``IPC'' and case names; age of the claim from the event date; language detection and a retry if an English user gets Hindi; \textbf{verdict rules} (amount over Rs 50 lakh, claim near its time limit, threats, a court case already filed, no matching law). \\ \hline
Human in the loop & The corpus is hand-reviewed with sources. The verdict recommends a lawyer, \emph{in addition to} steps, for the risky cases above. Drafts have \texttt{[PLACEHOLDERS]} and a ``check every detail'' note. \\ \hline
\end{tabularx}

\smallskip
\textbf{Swappable model:} one \texttt{llm()} function; set \texttt{LLM\_PROVIDER=openai} with a key, model and base URL to use any OpenAI-compatible API, or change \texttt{GEMINI\_MODEL}.

\section{How it avoids invented laws}
\begin{itemize}
\item \textbf{Closed corpus.} 53 entries (about 7k tokens) go in every prompt. The model may cite only by entry id; code drops any id not in the file and tells the user.
\item \textbf{Safety net on free text.} Section, rule or article numbers that appear nowhere in the file, ``IPC'', and ``X v.\ Y'' case citations are flagged to the user as unverified.
\item \textbf{Honest caveats.} Each entry says what is state-specific, a model law or unconfirmed (e.g.\ the Model Tenancy Act applies only where a state adopted it; RBI liability rules change for transactions from 1 Jan 2027; 3 entries are marked \emph{unverified}). Prompts forbid fixed refund deadlines, case law and certainty.
\item \textbf{Limit:} the model can still paraphrase an entry loosely. The citation check proves the entry exists, not that the summary fits.
\end{itemize}

\section{How I checked quality (real runs, honest failures)}
I wrote \textbf{17 test conversations} and ran them on the live site. Checks: verdict, issue type, a cited law id, reply language, forbidden text. Cases: refunds in English, Hindi and Telugu; three scam types; a Rs 60 lakh claim; a 30-month-old claim; salary; deposit; Telugu private-loan documents; cheque notice; blackmail; a trap asking for ``IPC 420 and a Supreme Court case''; a vague one-line message. \textbf{\textcolor{OliveGreen}{Final: 17/17 pass}}, but the path there is the useful part:

\smallskip
\begin{tabularx}{\linewidth}{|>{\raggedright\arraybackslash}p{33mm}|Y|}
\hline
\textbf{Run} & \textbf{What failed and the fix} \\ \hline
13 cases, run 1: 12/13 & The model said ``7 months'' for a May 2024 event (it does not know today's date), so a time-barred claim was called ``can help''. \textbf{Fix:} the model extracts the date, code computes the age. \\ \hline
After adding salary, deposit, loan, cheque: 14/17 & (a) The model set ``court case exists'' for a plain company notice, giving the wrong ``lawyer'' verdict. \textbf{Fix:} renamed the field to \texttt{court\_case\_filed} and tightened the prompt. (b) Gemini 503/429 overloads and a Vercel timeout: added retries. (c) Two failures were my test script giving up too early. \\ \hline
Earlier, by hand & English messages got Hindi or Hinglish replies (the prompt listed Hindi as an example). \textbf{Fix:} language detected in code, enforced, one retry if wrong. \\ \hline
\end{tabularx}

\smallskip
\textbf{Where it still fails or is weak:} (1) my cases test behaviour, not legal correctness; no lawyer has reviewed the output. (2) The cheque-bounce case gets ``can help'' although a bounced cheque carries criminal exposure; the verdict is too optimistic there. (3) Practice mode sometimes states a detail slightly wrong (said ``a whole month'' for about three weeks). (4) Deposit and wage law is state-specific, so answers are often general guidance. (5) Replies take 5--25\,s on the free tier, and a stalled Gemini call can exceed Vercel's 60\,s limit (one run failed this way). (6) Voice input depends on the browser; voice output uses the browser's voice.

\section{AI tools used, and where they got it wrong}
\textbf{Claude Code} wrote the code and helped draft the corpus; three Claude research sub-agents checked each legal claim against web sources (they were told to write NOT VERIFIED, not guess); \textbf{Gemini 3.5 Flash-Lite} runs inside the app. \textbf{Where they were wrong:} (1) Claude's recalled section numbers were wrong twice and the source check caught both: Sale of Goods Act \emph{s.56} is damages for \emph{non-acceptance} (non-delivery is \emph{s.57}), and Consumer Protection Act \emph{s.51} is an appeal, not State jurisdiction. This is why the corpus is built from sources, not memory. (2) Gemini mis-computed the age of a claim (above). (3) The default model \texttt{gemini-2.5-flash} had been retired for new users; larger Gemini models took 15--60\,s per call, so I chose a lite one.

\section{Scale and cost}
A case takes about \textbf{5--7 model calls}: intake 1--3 (small prompt), confirm and analysis, plan, then follow-ups. The corpus (about 7k tokens) is sent with analysis, plan, follow-up and practice calls, so roughly \textbf{25k input and 3k output tokens per case}. I have not measured rupee cost; it is those tokens times the chosen model's price. The server is stateless, so it scales horizontally and stores no personal data. \textbf{At volume:} a paid model (one env change); cache the fixed corpus prefix; once the corpus passes a few hundred entries, send only entries matching the detected issue type; queue and rate-limit calls; log anonymised failures for lawyer review; re-verify entries on a schedule (each has \texttt{verified\_on}) because law moves (RBI rules 2027, new labour codes, state rent laws).

\section{Next steps}
Have a lawyer review the corpus and 20 outputs; verify the 3 unverified and 33 secondary entries against official text; add a cheque-bounce ``criminal exposure'' rule; add state-specific rent and money-lender entries; measure real cost and latency; test with real users in Hindi, Telugu and Tamil.

\end{document}
